\ifdefined\gkver\else\def\gkver{student}\fi
\documentclass[\gkver]{gaokaozhenti}
\newcommand{\ccnum}[1]{{\fontspec{Noto Serif CJK SC}#1}}
\begin{document}

\chapter{2010年安徽理综}
%% paperType: 理综物理部分

\begin{enumerate}
\item
%% number: 14
%% paperName: 2010年安徽理综第 14 题 6 分
%% typeId: 1
%% type: 单选题
%% chapter: 功和能
%% point: 机械能守恒定律
%% method: 无
%% score: 6
%% degree: 650
%% duplicateId: 0
%% body:
伽利略曾设计如图所示的一个实验，将摆球拉至 M 点放开，摆球会达到同一水平高度上的 N 点。如果在 E 或 F 处钉钉子，摆球将沿不同的圆弧达到同一高度的对应点；反过来，如果让摆球从这些点下落，它同样会达到原水平高度上的 M 点。这个实验可以说明，物体由静止开始沿不同倾角的光滑斜面（或弧线）下滑时，其末速度的大小\xzanswer{C}

\onepicture[w=5.5cm]{figs/14.png}

\fourchoices[answer=C]
{只与斜面的倾角有关}
{只与斜面的长度有关}
{只与下滑的高度有关}
{只与物体的质量有关}

%% answer: C
%% memo:
\memoanswer{
伽利略的理想斜面和摆球实验，斜面上的小球和摆线上的小球好像“记得”起自己的起始高度，实质是动能与势能的转化过程中，总能量不变。物体由静止开始沿不同倾角的光滑斜面（或弧线）下滑时，高度越大，初始的势能越大转化后的末动能也就越大，速度越大。选项 C 正确。
}

\item
%% number: 15
%% paperName: 2010年安徽理综第 15 题 6 分
%% typeId: 1
%% type: 单选题
%% chapter: 机械波
%% point: 波的图象
%% method: 无
%% score: 6
%% degree: 650
%% duplicateId: 0
%% body:
一列沿 $x$ 轴方向传播的简谐横波，某时刻的波形如图所示。P 为介质中的一个质点，从该时刻开始的一段极短时间内，P 的速度 $v$ 和加速度 $a$ 的大小变化情况是\xzanswer{D}

\onepicture[w=5cm]{figs/15.png}

\fourchoices[answer=D]
{$v$ 变小，$a$ 变大}
{$v$ 变小，$a$ 变小}
{$v$ 变大，$a$ 变大}
{$v$ 变大，$a$ 变小}

%% answer: D
%% memo:
\memoanswer{
由题图可得，波沿 $x$ 轴方向传播，P 质点在该时刻的运动方向沿 $y$ 轴正方向，向平衡位置靠近，做加速度减小的加速运动，$v$ 变大，$a$ 变小，选项 D 正确。
}

\item
%% number: 16
%% paperName: 2010年安徽理综第 16 题 6 分
%% typeId: 1
%% type: 单选题
%% chapter: 电场
%% point: 电势差与电场强度的关系
%% method: 无
%% score: 6
%% degree: 650
%% duplicateId: 0
%% body:
如图所示，在 $xOy$ 平面内有一个以 O 为圆心、半径 $R=0.1\Um$ 的圆，P 为圆周上的一点，O、P 两点连线与 $x$ 轴正方向的夹角为 $\theta$。若空间存在沿 $y$ 轴负方向的匀强电场，场强大小 $E=100\UVm$，则 O、P 两点的电势差可表示为\xzanswer{A}

\onepicture[w=4.5cm]{figs/16.png}

\fourchoices[answer=A]
{$U_{\text{OP}}=-10\sin\theta$（V）}
{$U_{\text{OP}}=10\sin\theta$（V）}
{$U_{\text{OP}}=-10\cos\theta$（V）}
{$U_{\text{OP}}=10\cos\theta$（V）}

%% answer: A
%% memo:
\memoanswer{
在匀强电场中，两点间的电势差 $U=Ed$，而 $d$ 是沿场强方向上的距离，所以 $d_{\text{OP}}=-R\sin\theta$，故：$U_{\text{OP}}=-100\times0.1\sin\theta=-10\sin\theta$（V），选项 A 正确。
}

\item
%% number: 17
%% paperName: 2010年安徽理综第 17 题 6 分
%% typeId: 1
%% type: 单选题
%% chapter: 曲线运动
%% point: 人造地球卫星
%% method: 无
%% score: 6
%% degree: 450
%% duplicateId: 0
%% body:
为了对火星及其周围的空间环境进行探测，我国预计于 2011 年 10 月发射第一颗火星探测器“萤火一号”。假设探测器在离火星表面高度分别为 $h_{1}$ 和 $h_{2}$ 的圆轨道上运动时，周期分别为 $T_{1}$ 和 $T_{2}$。火星可视为质量分布均匀的球体，且忽略火星的自转影响，万有引力常量为 $G$。仅利用以上数据，可以计算出\xzanswer{A}

\fourchoices[answer=A]
{火星的密度和火星表面的重力加速度}
{火星的质量和火星对“萤火一号”的引力}
{火星的半径和“萤火一号”的质量}
{火星表面的重力加速度和火星对“萤火一号”的引力}

%% answer: A
%% memo:
\memoanswer{
由于万有引力提供探测器做圆周运动的向心力，则有

$G\frac{Mm}{(R+h_{1})^{2}}=m\left(\frac{2\pi}{T_{1}}\right)^{2}(R+h_{1})$；

$G\frac{Mm}{(R+h_{2})^{2}}=m\left(\frac{2\pi}{T_{2}}\right)^{2}(R+h_{2})$，可求得火星的质量 $M=\frac{4\pi^{2}(R+h_{1})^{3}}{GT_{1}^{2}}=\frac{4\pi^{2}(R+h_{2})^{3}}{GT_{2}^{2}}$ 和火星的半径

$R=\frac{\sqrt[3]{\left(\frac{T_{2}}{T_{1}}\right)^{2}h_{2}-h_{1}}}{1-\sqrt[3]{\left(\frac{T_{2}}{T_{1}}\right)^{2}}}$，根据密度公式得：$\rho=\frac{M}{V}=\frac{M}{\frac{4}{3}\pi R^{3}}=\frac{3M}{4\pi R^{3}}$。在火星表面的物体有 $G\frac{Mm}{R^{2}}=mg$，

可得火星表面的重力加速度 $g=\frac{GM}{R^{2}}$，故选项 A 正确。
}

\item
%% number: 18
%% paperName: 2010年安徽理综第 18 题 6 分
%% typeId: 1
%% type: 单选题
%% chapter: 电路
%% point: 动态分析
%% method: 无
%% score: 6
%% degree: 450
%% duplicateId: 0
%% body:
如图所示，M、N 是平行板电容器的两个极板，$R_{0}$ 为定值电阻，$R_{1}$、$R_{2}$ 为可调电阻，用绝缘细线将质量为 $m$、带正电的小球悬于电容器内部。闭合电键 S，小球静止时受到悬线的拉力为 $F$。调节 $R_{1}$、$R_{2}$，关于 $F$ 的大小判断正确的是\xzanswer{B}

\onepicture[w=4.4cm]{figs/18.png}

\fourchoices[answer=B]
{保持 $R_{1}$ 不变，缓慢增大 $R_{2}$ 时，$F$ 将变大}
{保持 $R_{1}$ 不变，缓慢增大 $R_{2}$ 时，$F$ 将变小}
{保持 $R_{2}$ 不变，缓慢增大 $R_{1}$ 时，$F$ 将变大}
{保持 $R_{2}$ 不变，缓慢增大 $R_{1}$ 时，$F$ 将变小}

%% answer: B
%% memo:
\memoanswer{
保持 $R_{1}$ 不变，缓慢增大 $R_{2}$ 时，由于 $R_{0}$ 和 $R_{2}$ 串联，$R_{0}$ 两端的电压减小，即平行板电容器的两个极板的电压 $U$ 减小，带电小球受到的电场力 $F_{\text{电}}=qE=q\frac{U}{d}$ 减小，悬线的拉力为 $F=\sqrt{(mg)^{2}+F_{\text{电}}^{2}}$ 将减小，选项 B 正确，A 错误。保持 $R_{2}$ 不变，缓慢增大 $R_{1}$ 时，$R_{0}$ 两端的电压不变，$F_{\text{电}}$ 不变，悬线的拉力为 $F$ 不变，C、D 错误。
}

\item
%% number: 19
%% paperName: 2010年安徽理综第 19 题 6 分
%% typeId: 1
%% type: 单选题
%% chapter: 力物体的平衡
%% point: 受力分析
%% method: 无
%% score: 6
%% degree: 450
%% duplicateId: 0
%% body:
L 型木板 P（上表面光滑）放在固定斜面上，轻质弹簧一端固定在木板上，另一端与置于木板上表面的滑块 Q 相连，如图所示。若 P、Q 一起沿斜面匀速下滑，不计空气阻力。则木板 P 的受力个数为\xzanswer{C}

\onepicture[w=4.5cm]{figs/19.png}

\fourchoices[answer=C]
{3}
{4}
{5}
{6}

%% answer: C
%% memo:
\memoanswer{
P、Q 一起沿斜面匀速下滑时，由于木板 P 上表面光滑，滑块 Q 受到重力、P 的支持力和弹簧沿斜面向上的弹力。木板 P 受到重力、斜面的支持力、斜面的摩擦力、Q 的压力和弹簧沿斜面向下的弹力，所以选项 C 正确。
}

\item
%% number: 20
%% paperName: 2010年安徽理综第 20 题 6 分
%% typeId: 1
%% type: 单选题
%% chapter: 电磁感应
%% point: 线圈过有界磁场
%% method: 等效替代
%% score: 6
%% degree: 250
%% duplicateId: 0
%% body:
如图所示，水平地面上方矩形区域内存在垂直纸面向里的匀强磁场，两个边长相等的单匝闭合正方形线圈Ⅰ和Ⅱ，分别用相同材料，不同粗细的导线绕制（Ⅰ为细导线）。两线圈在距磁场上界面 $h$ 高处由静止开始自由下落，再进入磁场，最后落到地面。运动过程中，线圈平面始终保持在竖直平面内且下边缘平行于磁场上边界。设线圈Ⅰ、Ⅱ落地时的速度大小分别为 $v_{1}$、$v_{2}$，在磁场中运动时产生的热量分别为 $Q_{1}$、$Q_{2}$。不计空气阻力，则\xzanswer{D}

\onepicture[w=4.2cm]{figs/20.png}

\fourchoices[answer=D]
{$v_{1}<v_{2}$，$Q_{1}<Q_{2}$}
{$v_{1}=v_{2}$，$Q_{1}=Q_{2}$}
{$v_{1}<v_{2}$，$Q_{1}>Q_{2}$}
{$v_{1}=v_{2}$，$Q_{1}<Q_{2}$}

%% answer: D
%% memo:
\memoanswer{
由于从同一高度下落，到达磁场边界时具有相同的速度 $v$，切割磁感线产生感应电流同时受到磁场的安培力 $F=\frac{B^{2}l^{2}v}{R}$，又 $R=\rho\frac{4l}{S}$（$\rho$ 为材料的电阻率，$l$ 为线圈的边长，$S$ 为单匝导线横截面积），所以安培力 $F=\frac{B^{2}lvS}{4\rho}$，此时加速度 $a=g-\frac{F}{m}$，且 $m=\rho_{0}S\cdot4l$（$\rho_{0}$ 为材料的密度），所以加速度 $a=g-\frac{B^{2}v}{16\rho\rho_{0}}$ 是定值，线圈Ⅰ和Ⅱ同步运动，落地速度相等 $v_{1}=v_{2}$。由能量守恒可得：$Q=mg(h+H)-\frac{1}{2}mv^{2}$，（$H$ 是磁场区域的高度），Ⅰ为细导线 $m$ 小，产生的热量小，所以 $Q_{1}<Q_{2}$。正确选项 D。
}

\item
%% number: 21
%% paperName: 2010年安徽理综第 21 题 3 分
%% typeId: 4
%% type: 实验题
%% chapter: 功和能
%% point: DIS研究机械能守恒定律
%% method: 无
%% score: 3
%% degree: 650
%% duplicateId: 0
%% body:
\begin{enumerate}
	\item
	\textbf{Ⅰ．}
	\begin{enumerate}
		\item
		在测定金属的电阻率实验中，用螺旋测微器测量金属丝的直径，示数如图 \ref{2010安徽21a} 所示，读数为\tkanswer[1.5cm]{0.617} $\Umm$。

		\item
		在用单摆测定重力加速度实验中，用游标为 20 分度的卡尺测量摆球的直径，示数如图 \ref{2010安徽21b} 所示，读数为\tkanswer[1.5cm]{0.675} $\Ucm$。
	\end{enumerate}

	\twopicture[numstyle=arabic, numA=1, numB=2, labelA=2010安徽21a, labelB=2010安徽21b, wA=3.2cm, wB=9cm, align=right]
	{figs/21a.png}
	{figs/21b.png}

	\item
	\textbf{Ⅱ．}太阳能是一种清洁、“绿色”能源。在我国上海举办的 2010 年世博会上，大量利用了太阳能电池。太阳能电池在有光照时，可以将光能转化为电能，在没有光照时，可以视为一个电学器件。某实验小组根据测绘小灯泡伏安特性曲线的实验方法，探究一个太阳能电池在没有光照时（没有储存电能）的 $I-U$ 特性。所用的器材包括：太阳能电池，电源 $E$，电流表 A，电压表 V，滑动变阻器 $R$，开关 S 及导线若干。
	\begin{enumerate}
		\item
		为了达到上述目的，请将图 \ref{2010安徽21c} 连成一个完整的实验电路图。

		\item
		该实验小组根据实验得到的数据，描点绘出了如图 \ref{2010安徽21d} 的 $I-U$ 图像。由图可知，当电压小于 $2.00\UV$ 时，太阳能电池的电阻\tkanswer[1.2cm]{很大}（填“很大”或“很小”）；当电压为 $2.80\UV$ 时，太阳能电池的电阻约为\tkanswer[1.5cm]{$1.0\times10^{3}$} $\UO$。
	\end{enumerate}

	\twopicture[numstyle=arabic, numA=1, numB=2, labelA=2010安徽21c, labelB=2010安徽21d, hA=4.6cm, hB=4.6cm, align=right]
	{figs/21c.png}
	{figs/21d.png}

	\item
	\textbf{Ⅲ．}利用图示装置进行验证机械能守恒定律的实验时，需要测量物体由静止开始自由下落到某点时的瞬时速度 $v$ 和下落高度 $h$。某班同学利用实验得到的纸带，设计了以下四种测量方案。

	（A）用刻度尺测出物体下落的高度 $h$，并测出下落时间 $t$，通过 $v=gt$ 计算出瞬时速度 $v_{0}$；

	（B）用刻度尺测出物体下落的高度 $h$，并通过 $v=\sqrt{2gh}$ 计算出瞬时速度；

	（C）根据做匀速直线运动时纸带上某点的瞬时速度，等于这点前后相邻两点间的平均速度，测算出瞬时速度，并通过计算出高度 $h$；

	（D）用刻度尺测出物体下落的高度 $h$，根据做匀速直线运动时纸带上某点的瞬时速度，等于这点前后相邻两点间的平均速度，测算出瞬时速度 $v_{0}$。

	以上方案中只有一种正确，正确的是\tkanswer[1.2cm]{D}。（填入相应的字母）

	\twopicture[showlabel=false, hA=5cm, hB=3.5cm, align=right]
	{figs/21e.png}
	{figs/21f.jpg}
\end{enumerate}

%% answer: 
\jdanswer{
\begin{enumerate}
	\item
	Ⅰ．\begin{enumerate}
		\item
		$0.617$（$0.616\sim0.619$）；
		\item
		$0.675$。
	\end{enumerate}
	\item
	Ⅱ．\begin{enumerate}
		\item
		电路连接如图所示；
		\item
		很大；$1.0\times10^{3}\UO$。
	\end{enumerate}
	\item
	Ⅲ．D
\end{enumerate}
}
%% memo:
\memoanswer{
（1）$0.5\Umm+11.7\times0.01\Umm=0.617\Umm$；

（2）$0.6\Ucm+15\times0.05\Umm=0.675\Ucm$。
}
\memoanswer{
（1）根据测绘小灯泡伏安特性曲线的实验方法，电路连接如图。

\onepicture[w=4.5cm]{figs/21c_an.jpg}

（2）在电压小于 $2.00\UV$ 时，由图可读出电流很小，由 $I=\frac{U}{R}$ 得 $R=\frac{U}{I}$，太阳能电池的电阻很大。

（3）当电压为 $2.80\UV$ 时，根据题图读出 $U$、$I$，由 $R=\frac{U}{I}$ 得：$R=1.0\times10^{3}\UO$。
}
\memoanswer{
物体由静止开始自由下落过程中受到空气阻力和纸带与打点计时器的摩擦阻力作用，不是自由落体运动，a、b 错误。物体下落的高度是用米尺测量的，不是计算的，c 错误。d 为验证机械能守恒定律的实验测量方案，正确。
}

\item
%% number: 22
%% paperName: 2010年安徽理综第 22 题 14 分
%% typeId: 6
%% type: 计算题
%% chapter: 运动和力
%% point: 牛顿定律的应用
%% method: 无
%% score: 14
%% degree: 450
%% duplicateId: 0
%% body:
质量为 $2\Ukg$ 的物体在水平推力 $F$ 的作用下沿水平面作直线运动，一段时间后撤去 $F$，其运动的 $v-t$ 图像如图所示。$g$ 取 $10\Umsq$，求：
\begin{enumerate}
	\item
	物体与水平面间的运动摩擦系数 $\mu$；

	\item
	水平推力 $F$ 的大小；

	\item
	$0\sim10\Us$ 内物体运动位移的大小。
\end{enumerate}

\onepicture[align=right, w=5cm]{figs/22.png}

%% answer: 
\jdanswer{
\begin{enumerate}
	\item
	$0.2$

	\item
	$6\UN$

	\item
	$46\Um$

\end{enumerate}
}
%% memo:
\memoanswer{
（1）设物体做匀减速运动的时间为 $\Delta t_{2}$、初速度为 $v_{20}$、末速度为 $v_{2t}$、加速度为 $a_{2}$，则

$a_{2}=\frac{v_{2t}-v_{20}}{\Delta t_{2}}=-2\Umsq$ ①

设物体所受的摩擦力为 $F_{f}$，根据牛顿第二定律有

$F_{f}=ma_{2}$ ②

$F_{f}=-\mu mg$ ③

联立②③得：$\mu=\frac{-a_{2}}{g}=0.2$ ④

（2）设物体做匀加速运动的时间为 $\Delta t_{1}$、初速度为 $v_{10}$、末速度为 $v_{1t}$、加速度为 $a_{1}$，则

$a_{1}=\frac{v_{1t}-v_{10}}{\Delta t_{1}}=1\Umsq$ ⑤

根据牛顿第二定律有 $F+F_{f}=ma_{1}$ ⑥

联立③⑥得：$F=\mu mg+ma_{1}=6\UN$

（3）解法一：由匀变速运动的位移公式得：

$x=x_{1}+x_{2}=v_{10}\Delta t_{1}+\frac{1}{2}a_{1}\Delta t_{1}^{2}+v_{20}\Delta t_{2}+\frac{1}{2}a_{2}\Delta t_{2}^{2}=46\Um$

解法二：根据 $v-t$ 图像围成的面积得：

$x=\left(\frac{v_{10}+v_{1t}}{2}\times\Delta t_{1}+\frac{1}{2}\times v_{20}\times\Delta t_{2}\right)=46\Um$
}

\item
%% number: 23
%% paperName: 2010年安徽理综第 23 题 16 分
%% typeId: 6
%% type: 计算题
%% chapter: 磁场
%% point: 带电粒子在磁场中的运动
%% method: 无
%% score: 16
%% degree: 450
%% duplicateId: 0
%% body:
如图 \ref{2010安徽23a} 所示，宽度为 $d$ 的竖直狭长区域内（边界为 $L_{1}$、$L_{2}$），存在垂直纸面向里的匀强磁场和竖直方向上的周期性变化的电场（如图 \ref{2010安徽23b} 所示），电场强度的大小为 $E_{0}$，$E>0$ 表示电场方向竖直向上。$t=0$ 时，一带正电、质量为 $m$ 的微粒从左边界上的 $N_{1}$ 点以水平速度 $v$ 射入该区域，沿直线运动到 Q 点后，做一次完整的圆周运动，再沿直线运动到右边界上的 $N_{2}$ 点。Q 为线段 $N_{1}N_{2}$ 的中点，重力加速度为 $g$。上述 $d$、$E_{0}$、$m$、$v$、$g$ 为已知量。
\begin{enumerate}
	\item
	求微粒所带电荷量 $q$ 和磁感应强度 $B$ 的大小；

	\item
	求电场变化的周期 $T$；

	\item
	改变宽度 $d$，使微粒仍能按上述运动过程通过相应宽度的区域，求 $T$ 的最小值。
\end{enumerate}

\twopicture[numstyle=arabic, numA=1, numB=2, labelA=2010安徽23a, labelB=2010安徽23b, hA=4cm, hB=4cm, align=right]
{figs/23a.png}
{figs/23b.png}

%% answer: 
\jdanswer{
\begin{enumerate}
	\item
	$q=\frac{mg}{E_{0}}$，$B=\frac{2E_{0}}{v}$

	\item
	$T=\frac{d}{2v}+\frac{\pi v}{g}$

	\item
	$T_{\min}=\frac{(2\pi+1)v}{2g}$

\end{enumerate}
}
%% memo:
\memoanswer{
（1）微粒做直线运动，则

$mg+qE_{0}=qvB$ ①

微粒做圆周运动，则

$mg=qE_{0}$ ②

联立①②得：$q=\frac{mg}{E_{0}}$ ③

$B=\frac{2E_{0}}{v}$ ④

（2）设微粒从 $N_{1}$ 运动到 Q 的时间为 $t_{1}$，做圆周运动的周期为 $t_{2}$，则

$\frac{d}{2}=vt_{1}$ ⑤

$qvB=m\frac{v^{2}}{R}$ ⑥

$2\pi R=vt_{2}$ ⑦

联立③④⑤⑥⑦得：$t_{1}=\frac{d}{2v}$，$t_{2}=\frac{\pi v}{g}$ ⑧

电场变化的周期 $T=t_{1}+t_{2}=\frac{d}{2v}+\frac{\pi v}{g}$ ⑨

（3）若微粒能完成题述的运动过程，要求 $d\geq2R$ \ccnum{⑩}

联立③④⑥得：$R=\frac{v^{2}}{2g}$ \ccnum{⑪}

设 $N_{1}Q$ 段直线运动的最短时间 $t_{1\min}$，由⑤\ccnum{⑩⑪}得

$t_{1\min}=\frac{v}{2g}$

因 $t_{2}$ 不变，$T$ 的最小值 $T_{\min}=t_{1\min}+t_{2}=\frac{(2\pi+1)v}{2g}$
}

\item
%% number: 24
%% paperName: 2010年安徽理综第 24 题 20 分
%% typeId: 6
%% type: 计算题
%% chapter: 运动和力
%% point: 动量和能量
%% method: 无
%% score: 20
%% degree: 250
%% duplicateId: 0
%% body:
如图，ABD 为竖直平面内的光滑绝缘轨道，其中 AB 段是水平的，BD 段为半径 $R=0.2\Um$ 的半圆，两段轨道相切于 B 点，整个轨道处在竖直向下的匀强电场中，场强大小 $E=5.0\times10^{3}\UVm$。一不带电的绝缘小球甲，以速度 $v_{0}$ 沿水平轨道向右运动，与静止在 B 点带正电的小球乙发生弹性碰撞。已知甲、乙两球的质量均为 $m=1.0\times10^{-2}\Ukg$，乙所带电荷量 $q=2.0\times10^{-5}\UC$，$g$ 取 $10\Umsq$。（水平轨道足够长，甲、乙两球可视为质点，整个运动过程无电荷转移）
\begin{enumerate}
	\item
	甲乙两球碰撞后，乙恰能通过轨道的最高点 D，求乙在轨道上的首次落点到 B 点的距离；

	\item
	在满足（1）的条件下，求甲的速度 $v_{0}$；

	\item
	若甲仍以速度 $v_{0}$ 向右运动，增大甲的质量，保持乙的质量不变，求乙在轨道上的首次落点到 B 点的距离范围。
\end{enumerate}

\onepicture[align=right, w=7cm]{figs/24.png}

%% answer: 
\jdanswer{
\begin{enumerate}
	\item
	$x=0.4\Um$

	\item
	$v_{0}=2\sqrt{5}\Ums$

	\item
	$0.4\Um\leq x'<1.6\Um$

\end{enumerate}
}
%% memo:
\memoanswer{
（1）在乙恰能通过轨道的最高点的情况下，设乙到达最高点的速度为 $v_{D}$，乙离开 D 点到达水平轨道的时间为 $t$，乙的落点到 B 点的距离为 $x$，则

$m\frac{v_{D}^{2}}{R}=mg+qE$ ①

$2R=\frac{1}{2}\left(\frac{mg+qE}{m}\right)t^{2}$ ②

$x=v_{D}t$ ③

联立①②③得：$x=0.4\Um$ ④

（2）设碰撞后甲、乙的速度分别为 $v_{\text{甲}}$、$v_{\text{乙}}$，根据动量守恒和机械能守恒定律有：

$mv_{0}=mv_{\text{甲}}+mv_{\text{乙}}$ ⑤

$\frac{1}{2}mv_{0}^{2}=\frac{1}{2}mv_{\text{甲}}^{2}+\frac{1}{2}mv_{\text{乙}}^{2}$ ⑥

联立⑤⑥得：$v_{\text{乙}}=v_{0}$ ⑦

由动能定理得：$-mg\cdot2R-qE\cdot2R=\frac{1}{2}mv_{D}^{2}-\frac{1}{2}mv_{\text{乙}}^{2}$ ⑧

联立①⑦⑧得：$v_{0}=\sqrt{\frac{5(mg+qE)R}{m}}=2\sqrt{5}\Ums$ ⑨

（3）设甲的质量为 $M$，碰撞后甲、乙的速度分别为 $v_{M}$、$v_{m}$，根据动量守恒和机械能守恒定律有：

$Mv_{0}=Mv_{M}+mv_{m}$ \ccnum{⑩}

$\frac{1}{2}Mv_{0}^{2}=\frac{1}{2}Mv_{M}^{2}+\frac{1}{2}mv_{m}^{2}$ \ccnum{⑪}

联立\ccnum{⑩⑪}得：$v_{m}=\frac{2Mv_{0}}{M+m}$ \ccnum{⑫}

由\ccnum{⑫}和 $M\gg m$，可得：$v_{D}\leq v_{m}<2v_{D}$ \ccnum{⑬}

设乙球过 D 点的速度为 $v_{D}'$，由动能定理得

$-mg\cdot2R-qE\cdot2R=\frac{1}{2}mv_{D}'^{2}-\frac{1}{2}mv_{m}^{2}$ \ccnum{⑭}

联立⑨\ccnum{⑬⑭}得：$2\Ums\leq v_{D}'<8\Ums$ \ccnum{⑮}

设乙在水平轨道上的落点到 B 点的距离为 $x'$，则有：

$x'=v_{D}'t$ \ccnum{⑯}

联立②\ccnum{⑮⑯}得：$0.4\Um\leq x'<1.6\Um$
}

\end{enumerate}

\end{document}
